\subsection{向量丛的定义}

7.1.1. 一个向量丛图 \(t = (U, \varphi, F)\) 是 M 的一个三元组，其中 U 是 B 的开子集，F 是巴拿赫空间，而 \(\varphi\) 是从 \(\pi^{-1}(U)\) 到 \(U \times F\) 的双射，满足 \(\pi(\varphi^{-1}(b, h)) = b\)，对所有 \(b \in B\) 和所有 \(h \in F\) 均成立。我们称 U 为向量丛图 \(t\) 的定义域，并称 \(t\) 是 M 在 \(b \in B\) 处的向量丛图，当 \(b \in U\) 时。对每个 \(b \in U\)，记 \(t_b\) 为从 F 到 \(M_b\) 的双射，其定义为 \(t_b(h) = \varphi^{-1}(b, h)\)，其中 \(h \in F\)。

7.1.2. 设 \(t = (\mathrm{U}, \varphi, \mathrm{F})\) 和 \(t' = (\mathrm{U}', \varphi', \mathrm{F}')\) 是 \(\mathbf{M}\) 的两个向量丛图。它们称为 \(\mathbf{C}'\)-相容（或简称相容），当且仅当存在一个映射 \(\lambda\)，其类为 \(\mathbf{C}'\)，从流形 \(\mathrm{U} \cap \mathrm{U}'\) 到巴拿赫空间 \(\mathcal{L}(\mathrm{F}; \mathrm{F}')\)，使得：

\[
t _ {b} = t _ {b} ^ {\prime} \circ \lambda (b) \quad \text { for all } b \in \mathbf {U} \cap \mathbf {U} ^ {\prime}.
\]

7.1.3. M 的一组向量丛图称为 M 的 C'-向量丛图册（或简称向量丛图册），如果它由两两 C'-相容且定义域的并为 B 的向量丛图组成。M 的两个向量丛图册 A 和 B 称为 C'-等价（或等价），如果 A ∪ B 仍是 M 的向量丛图册。这个关系是一个等价关系。

7.1.4. M 上的类为 \(C^{r}\) 的向量丛结构（基底 B）是一个等价向量丛图册类的数据（《集合论》，第 II 章，第 6 节，第 9 号）。属于这个类中某个向量丛图册的向量丛图称为向量丛 M 的向量丛图。

设 M 是以 B 为基底的向量丛。对每个 \(b \in B\)，纤维 \(M_{b}\) 上存在唯一一个巴拿赫空间结构，使得对于向量丛 M 在 b 处的每个向量丛图 \(t = (\mathbf{U}, \varphi, \mathbf{F})\)，映射 \(t_{b}\) 都是从 F 到 \(M_{b}\) 的同构。

设 \(c = (\mathbf{U}, \psi, \mathbf{E})\) 是流形 \(\mathbf{B}\) 的一个图，\(t = (\mathbf{U}, \varphi, \mathbf{F})\) 是向量丛 \(\mathbf{M}\) 的一个向量丛图，二者有相同的定义域 \(\mathbf{U}\)。对 \(x \in \pi^{-1}(\mathbf{U})\)，置：

\[
\alpha (x) = (\psi (\pi (x)), t _ {\pi (x)} ^ {- 1} (x)).
\]

于是，三元组 \((\pi^{-1}(\mathbf{U}), \alpha, \mathbf{E} \times \mathbf{F})\) 是集合 M 的一个图。M 上存在唯一一个类为 \(C^{r}\) 的流形结构（称为 M 的底层结构），使得如此得到的所有图都是流形 M 的图。于是三元组 \((\mathbf{M}, \mathbf{B}, \pi)\) 是一个纤维化（6.1.1）。

7.1.5. 设 F 是一个巴拿赫空间。置 \(M = B \times F\)，其中映射取为第一投影。M 上存在唯一一个向量丛结构（以 B 为基底），使得 \((\mathbf{B}, \mathbf{Id}_{\mathbf{M}}, \mathbf{F})\) 是一个向量丛图。我们称赋予此结构的 \(B \times F\) 为以 B 为基底、纤维为 F 的平凡向量丛，有时记为 \(F_{B}\)。\(F_{B}\) 的流形结构是乘积流形结构，并且对于每个 \(b \in B\)，映射 \(h \mapsto (b, h)\) 是从 F 到 \(F_{B}\) 的纤维的巴拿赫空间同构，该纤维位于点 \(b \in B\) 处。纤维退化为 0 的平凡向量丛通常记为 0。

7.1.6. 设 M 是以 B 为基底的向量丛。对于 \(b \in B\)，将（有限或 \(+\infty\) 的）巴拿赫空间 \(M_{b}\) 的维数称为 M 在 b 处的秩，记为 \(\operatorname{rg}_{b}(\mathbf{M})\)。有 \(\dim_{x}\mathbf{M} = \dim_{b}\mathbf{B} + \operatorname{rg}_{b}\mathbf{M}\)，其中 \(b = \pi(x)\)。函数 \(b \mapsto \operatorname{rg}_{b}\mathbf{M}\) 局部常值。若满足 \(\operatorname{rg}_{b}\mathbf{M} < +\infty\) 对所有 \(b \in B\)，则称 M 是有限秩的。

